\section{FF-PGRPO Analyses: Feasibility Before Positive Credit}
\label{sec:ffpgrpo-analysis}

Aggregate reward can hide an unsafe trade: progress or comfort may increase
enough to give positive standardized advantage even when collision avoidance,
drivable-area compliance, or a protected reference metric regresses.  The
purpose of FF-PGRPO is therefore not merely to reshape the scalar score, but to
change which rollouts are eligible for positive credit.

\subsection{What the stage ablation establishes}
\label{sec:ffpgrpo-stage-ablation}

On NAVSIM v1, the no-stage baseline in the main ablation has 86.5 PDMS, 98.1
NC, 94.7 DAC, 94.2 TTC, and 80.9 EP.  Applying FF-PGRPO without PC-MTS reaches
91.0 PDMS, with 98.4 NC, 98.0 DAC, 95.8 TTC, and 85.93 EP.  With PC-MTS as the
initialization, FF-PGRPO reaches 91.1 PDMS, 98.5 NC, 98.0 DAC, 96.0 TTC, and
86.0 EP before APR.  The improvement is therefore not produced by sacrificing
the displayed hard-safety components for progress: DAC improves by 3.3 points
and TTC by 1.6 points in the FF-PGRPO-only row, while EP increases by 5.03
points.  These are paper-reported point estimates whose run/seed aggregation
is not stated, so the small 0.1-point
difference between the two FF-PGRPO rows is not presented as a stable gain.

\begin{table*}[t]
\centering
\small
\caption{Stage-wise NAVSIM v1 ablation under single-trajectory inference. Higher is better. Rows are one paper-reported run; only the final AMPT row is independently reproduced from 12,138 scene scores.}
\label{tab:stage_ablation}
\resizebox{\textwidth}{!}{%
\begin{tabular}{lcccccccccccc}
\toprule
configuration & pcmts & scalar\_grpo & ff\_pgrpo & apr & NC & DAC & TTC & C & EP & PDMS & evidence & runs \\
\midrule
GT-only & no & no & no & no & 98.1 & 94.7 & 94.2 & 100.0 & 80.9 & 86.5 & paper reported & 1 reported \\
PC-MTS & yes & no & no & no & 98.2 & 95.2 & 94.5 & 100.0 & 81.3 & 86.9 & paper reported & 1 reported \\
Scalar optimization & no & yes & no & no & 98.4 & 98.0 & 95.8 & 100.0 & 85.93 & 91.0 & paper reported & 1 reported \\
PC-MTS + FF-PGRPO & yes & no & yes & no & 98.5 & 98.0 & 96.0 & 100.0 & 86.0 & 91.1 & paper reported & 1 reported \\
Full AMPT & yes & no & yes & yes & 98.5 & 98.0 & 96.0 & 100.0 & 86.7 & 91.4 & paper reported / final row scene-reproduced & 1 reported + 12138 scenes final \\
\bottomrule
\end{tabular}%
}
\end{table*}


This result is consistent with the proposed ordering---recover feasibility,
then optimize continuous quality---but an aggregate stage ablation cannot
directly reveal the sign of individual rollout credits.  The fixed hard-scene
transition in Section~\ref{sec:hard658} provides scene-level evidence, and the
instrumentation below defines the rollout-level test.

\subsection{Credit diagnostics}
\label{sec:ffpgrpo-credit-diagnostics}

For rollout $i$, let $A_i>0$ denote positive assigned advantage and let
$V_i=1$ denote a violation of any hard-feasibility, protected-reference, or
admissible-Pareto condition.  We report
\begin{equation}
  \mathrm{PCVR} =
  \frac{\sum_i \mathbb{1}[A_i>0]V_i}
       {\sum_i \mathbb{1}[A_i>0]},
  \label{eq:positive-credit-violation}
\end{equation}
and separately count reference-regressing and Pareto-dominated positive
credits.  A denominator of zero is reported as ``no positive credits,'' not as
zero violation.  FF-PGRPO's implementation makes a violating rollout
ineligible for positive credit and exposes assertions for this contract.  The
archived training run, however, does not contain the per-rollout trace needed
to quote an observed PCVR; we do not convert a code invariant into a fabricated
measured rate.

The stepwise mechanism ablation adds: (i) the hard feasibility gate, (ii) the
coherent-reference guard, (iii) the Pareto positive-credit gate, and (iv) the
all-infeasible recovery rule.  Every variant logs group state
(all-feasible/mixed/all-infeasible), feasibility rate, positive and negative
advantage means, PCVR, reference-regressing positive-credit rate,
dominated-positive-credit rate, zero-to-feasible repair, positive-to-zero
regression, and updates until the first feasible rollout.  The coherent
reference is always one executable trajectory.  The historical
component-wise envelope remains only as an ablation because combining the best
value of each metric from different trajectories may describe no realizable
plan and can make every sampled rollout appear to regress.

The low-cost diagnostic run uses the paper's group size of 16 and stops at 300
updates; the final run uses the reported 10 epochs.  All variants share the
same sampled rollout tensors when credit is recomputed offline, which makes
the initial credit-assignment comparison independent of policy drift.  New
training is required only for the subsequent learning curves and remains
disabled unless explicitly enabled.

\subsection{Hard-scene evidence for feasibility-first recovery}
\label{sec:ffpgrpo-hard-evidence}

The manuscript recovery checkpoint is evaluated on the same 658-scene set as
scalar GRPO.  Scalar GRPO produces a positive score in 367 scenes (55.8\%,
95\% Wilson CI 52.0--59.5\%), whereas the AMPT recovery checkpoint is positive
in 440 scenes (66.9\%, CI 63.2--70.4\%).  The comparison uses a single
trajectory per scene and neither method uses test-time sampling, scoring, or
reranking.  Crucially, the paired transition is not merely a difference of
gross counts: 93 scalar-GRPO zeros become positive, 20 scalar-GRPO positives
become zero, 347 successes are retained, and 198 failures persist.  Net
recovery is therefore $93-20=73$ scenes, or 11.1 percentage points.  A
10,000-resample scene-paired bootstrap (analysis seed 2027) gives a 95\% CI of
8.1--14.3 percentage points for this paired change.  This interval quantifies
scene variation for this fixed checkpoint pair; it does not replace missing
multi-seed training uncertainty.

The subset mean PDMS increases from 0.5101 to 0.6145.  The corresponding
component means change from 0.8381 to 0.8815 for NC, 0.7052 to 0.7781 for DAC,
0.7705 to 0.8161 for TTC, and 0.4956 to 0.5932 for EP; comfort increases from
0.9985 to 1.0000, while DDC changes from 0.9567 to 0.9521.  This last decrease
is why protected metrics must be reported next to aggregate recovery rather
than summarized as universal component-wise improvement.  The complete
transition and failure-type decomposition are given in
Section~\ref{sec:hard658}.

The aggregate ablation, paired hard-scene transition, and positive-credit
instrumentation serve complementary roles.  The first shows the overall stage
gain, the second verifies that recovery is net positive rather than purchased
entirely with new failures, and the third is the predeclared test of the
claimed credit-assignment mechanism.
